% Chapter 11 -- company card. Update with companies.py, web/identity.py, web/profile.py,
% web/people.py, web/products.py, web/screening.py, web/social.py, registry/*, financials/* (trend.py, paste.py,
% info_sites.py), gui/company_card.py, gui/social_links.py, gui/financials_paste.py, gui/companies_panel.py,
% gui/affinity_box.py, applications/detail_widgets.py, jobsearch/employers.py (seed of a card from a posting).
\chapter{Company card}\label{ch:card}

A card is created when you \textbf{star} a company in Results (it becomes a \emph{favourite}), save it by hand,
import it, or press \ui{Create company card} on a posting (chapter~\ref{ch:jobs}). There is one card per company; a
new search updates the search data on it and \emph{never} touches your own fields (status: interesting / applied /
interview / offer / closed, notes, contacts, tags). Saving a result never uses the network: everything below runs
only when you ask. Above the sections the card shows the \textbf{affinity box} (chapter~\ref{ch:affinity}); below
them the box \emph{Applications to this company} (chapter~\ref{ch:tracker}).

\section{The nine sections}
{\small
\begin{tabularx}{\textwidth}{@{}l Y@{}}
\toprule
\hdr \thead{Section} & \thead{Content and origin} \\
\midrule
1. Header          & Name, legal name, REA, ATECO, tax code, VAT number. \\
2. Registry data   & Legal form, registered office, sites, PEC, phone, e-mail, share capital, founding year, size hints, LEI, own website. \\
3. Profiles sought & Careers snapshot: postings with attributes, open-application channel, date checked. \\
4. Products and services & From the home page, with source URLs. \\
5. Selection system & ATS and what the \emph{platform} offers, per its vendor. \\
6. CV screening    & Explicit statements on the employer's own pages (automated decisions, AI), quoted with URL. \\
7. Presentation and history & Sentences from the company's own pages. \\
8. People and contacts & Role mailboxes and named people, HR first. \\
9. Financials      & Revenue and net profit per year, chart and revenue trend; always shown, with the reason when empty. \\
\bottomrule
\end{tabularx}
}

\section{Links of the company}\label{sec:links}
A row of buttons beside the name opens the website and the profiles, in this order when known: Website, LinkedIn,
Facebook, Instagram, X, YouTube, TikTok, GitHub, XING, Glassdoor, Indeed, Wikipedia, Wikidata, Google Maps (also
\ui{Links of this company} in the Results and Companies menus). \textbf{A profile address is never guessed}: it
appears only if the company's own site, OpenStreetMap, Wikidata or the card names it; share and personal links are
rejected. With no LinkedIn page known, \ui{LinkedIn (search)} is offered, and \ui{Google Maps (search)} always,
drawn dashed: they only open a search in \emph{your} browser. Every tooltip says where the link was found.

\section{Building a card}
\ui{Create full card} (\ui{Refresh full card} once built), \ui{Stop} and \ui{More $>$ Refresh section} run in the
background; every step is \emph{saved when done} and \ui{Stop} keeps what was read. Steps: website; legal, contacts,
about and careers pages; registers (VIES, GLEIF, Wikidata); presentation and history; people; financial figures;
selection system and privacy notices; products. All is polite (shared fetcher, at most one VIES or GLEIF request per
2\,s per host, answers cached 30 days). \textbf{Batch:} \ui{Full cards for selected} and \ui{Full cards for all
favourites} build one card after the other with a small ``3/12'' bar; a failure is logged and skipped.

\section{Identity from the website, and the cross-check}
Legal identifiers are read from the company's \emph{own} site: the home page plus up to six pages chosen by link text
and path (legal notice, contacts, about, history, careers, team, privacy). Extracted: VAT number and tax code (check
digit verified, a wrong one discarded), REA, share capital, PEC, e-mail, phone, registered office, sites, founding
year, size hints, schema.org data. Every value keeps its \emph{provenance}: page URL, place on the page, date. The
identifiers are compared with \textbf{VIES} (``not valid'' only means ``not enabled for intra-EU trade''),
\textbf{GLEIF} and \textbf{Wikidata}, searched \emph{only by an identifier}, never by name; register data you enter or
import is a further source.

{\small
\begin{tabularx}{\textwidth}{@{}l Y@{}}
\toprule
\hdr \thead{Status} & \thead{When} \\
\midrule
\textcolor{othGreen}{\textbf{CONFIRMED}}  & A checksum-valid VAT/tax code on the site equals one in an independent register, the names agree, and the town when the site states it. \\
\textcolor{othBlue}{\textbf{PROBABLE}}    & Valid site identifier but no register hit, or a register match by name only, or a candidate that failed one check. \\
\textcolor{othGray}{\textbf{UNVERIFIED}}  & No identifier found. \\
\textcolor{othRed}{\textbf{CONFLICT}}     & The same kind of identifier differs between sources. \\
\bottomrule
\end{tabularx}
}

\vspace{2pt}
\ui{Why?} shows ``Agreement between sources''. When sources give \emph{different values} for REA, share capital, PEC,
founding year, legal form, registered office or LEI, all are kept and shown as a conflict. In section 5 the link
\ui{Simulate this system on my CV} opens tab \ui{Recruiter analysis} with this company and its system preselected
(chapter~\ref{ch:analysis}).

\section{Presentation, history, people, blanks}
\textbf{Presentation and history} are sentences copied from the company's own pages, cut at a sentence end, never
rewritten. \textbf{People and contacts}: role mailboxes sorted into \emph{HR} (first), \emph{management},
\emph{general} and \emph{privacy} (never proposed for an application); named people only when a page states name
and role together; the chief executive from Wikidata. An address is \emph{never guessed}; contacts stay on this
computer and can be deleted. \textbf{Blanks stay blank}: a value not found shows ``---'' with a hint on where to
get it free by hand. \ui{More $>$ Open Registro Imprese} opens the free search and copies the best query;
\ui{Enter data from Registro Imprese\ldots} takes ATECO, PEC, REA, tax code, legal form and registered office,
checked first (checksums, REA format, ATECO against the 2025 table) and stored as ``Registro Imprese (entered by
hand)'' with the date. With a provider and \emph{your own key}, \ui{Fetch from paid register} makes \emph{one}
lookup after showing the cost and asking you to confirm; no build or batch ever uses it.

\begin{warnbox}
The paid-provider adapter was written from the provider's documentation and tested only with invented answers: it has
\emph{not} been tried against the live service. Check the first result by hand.
\end{warnbox}

\section{Financials}
Section 9 \emph{is always shown}. With no figures it says why (``Not checked yet: use Create full card'', or the
reason of the last attempt) and where the numbers are: the balance sheet a non-listed Italian company files with the
Registro Imprese (about 2.50\,EUR a year there; public sites often show ranges). With figures it charts revenue and
net profit per fiscal year (a loss is hatched) and employees, with source, year, currency and scope for each.
\textbf{Free figures exist only for some companies:} listed companies (ESEF reports on \code{filings.xbrl.org},
by LEI, consolidated) and, rarely, Wikidata. Otherwise: \ui{Open company-information sites} (a web search the
program never reads), \ui{Paste figures\ldots}, \ui{Import balance-sheet CSV\ldots} (an ambiguous ``1.234'' is
rejected, not guessed) or \ui{Add figures by hand\ldots}. Hand-typed figures win over ESEF, then the paid provider,
then Wikidata; a disagreeing source is kept and flagged; \emph{consolidated} and \emph{separate} figures are never
mixed; no currency is converted.

\textbf{Paste figures.} Copy the figures from a page, paste them and press \ui{Read the figures}: a grid (Year,
Revenue, Net profit, Employees) appears for you to correct. Understood: Italian and English labels (\emph{Fatturato,
Ricavi, Revenue, Turnover}; \emph{Utile, Perdita}; \emph{Dipendenti, Addetti}), one figure per line or a table with
years in its header, amounts such as ``1.234.567'', ``1,2 mln'', ``850 mila'', ``(120)'' as negative; a \emph{range}
(``10--50 milioni'') stays a range. A number followed by \% is a rate, never an amount, and is ignored. The grid shows amounts with the thousands separators of the interface language (``12,345,678'' in English). A source text is required, nothing is saved until you confirm, and figures are
stored as ``pasted and reviewed''.

\textbf{Revenue trend.}\label{sec:trend} The rule is the yearly growth rate (CAGR) of revenue between the first and
the last year with revenue, for one scope and currency:
$\bigl(R_{\mathrm{last}}/R_{\mathrm{first}}\bigr)^{1/(\mathrm{last}-\mathrm{first})}-1$ over the latest ten years,
\emph{only with at least three yearly values} and both ends above zero; otherwise ``not enough data (at least 3 years
needed)''. From $+2$\,\%/year the revenue is \emph{growing}, at $-2$\,\% or less \emph{declining}, in between
\emph{stable}. It appears as a badge and as the \ui{Revenue trend} column, from stored figures only; a growing revenue
earns the small ranking bonus of chapter~\ref{ch:discovery}.

\section{Favourites and export}
The star in Results, the menu-bar star and \ui{File $>$ Favourites} (\code{Ctrl+Shift+F}) open the favourites window.
\ui{My companies} has filters, search by meaning, similar companies and hand-added companies. \ui{File $>$ Export CSV}
writes the Results table plus site, careers page, postings, workplace, origin, travel, spontaneous, blocked, favourite
and \emph{why selected}. A card can be copied as text, exported as Markdown, or printed / saved with a preview
(chapter~\ref{ch:apply}).
